% Chapter 4: Methodology (proposed approach)
% Rename sections to suit your work, e.g. "System Design" or "Proposed Model".
\chapter{Methodology}\label{ch:methodology}

\guide{Describe what you built or did, in enough detail that someone else
could reproduce it. Justify each design choice and relate it back to the
research questions.}

Give a short introduction to the chapter here.

%----------------------------------------------------------------------------
\section{Research Design}\label{sec:design}

Describe the overall research approach and why it suits the research
questions.

%----------------------------------------------------------------------------
\section{Dataset}\label{sec:dataset}

\subsection{Data Collection}
Describe where the data comes from, its size and how it is split.

\subsection{Preprocessing}
Describe how the data is cleaned and prepared.

%----------------------------------------------------------------------------
\section{Proposed Architecture}\label{sec:architecture}

Every figure must be referred to in the text. \Cref{fig:architecture} shows
the overall architecture of the proposed approach.

\begin{figure}[htbp]
  \centering
  % Replace example-image with your own file in figures/,
  % e.g. \includegraphics[width=0.8\textwidth]{architecture.pdf}
  \includegraphics[width=0.7\textwidth]{example-image}
  \caption{Overall architecture of the proposed approach}\label{fig:architecture}
\end{figure}

\subsection{[Component One]}
Describe the first component.

\subsection{[Component Two]}
Describe the second component. \Cref{alg:example} shows how an algorithm is
written.

\begin{algorithm}[htbp]
  \caption{Example algorithm}\label{alg:example}
  \begin{algorithmic}[1]
    \Require input $x$
    \Ensure output $y$
    \State $y \gets 0$
    \For{$i \gets 1$ \textbf{to} $n$}
      \State $y \gets y + f(x_i)$
    \EndFor
    \State \Return $y$
  \end{algorithmic}
\end{algorithm}

%----------------------------------------------------------------------------
\section{Training Procedure}\label{sec:training}

Describe the loss function, optimiser and training strategy. Number equations
so they can be referred to, e.g.\ \cref{eq:loss}:
\begin{equation}
  \mathcal{L} = \frac{1}{N} \sum_{i=1}^{N} \ell\left(\hat{y}_i, y_i\right).
  \label{eq:loss}
\end{equation}

%----------------------------------------------------------------------------
\section{Evaluation Metrics}\label{sec:evaluation-metrics}

Define the metrics used to evaluate the approach and explain why they were
chosen.

%----------------------------------------------------------------------------
\section{Ethical Considerations}\label{sec:ethics}

Discuss any ethical issues, for example the privacy of people in the data.
